\section{Introduction}
A \emph{module-finite extension} $R\subseteq S$ is an inclusion of commutative rings with the same identity such that $S$ is finitely generated as an $R$-module. It splits over $R$ if some $R$-linear functional $\rho:S\to R$ takes $1_S$ to $1$. All rings are Noetherian. The direct summand theorem settles the question over a regular base \cite{Andre}. For a singular local ring $(R,\m,k)$, Koh asked whether finite projective dimension of the extension suffices \cite{Koh}. We say that $R$ has the \emph{Koh property}, denoted by $\K$, if the following holds.
\[
 \text{Every module-finite extension }R\subseteq S
 \text{ with }\pd_R S<\infty\text{ splits over }R.
\]
We allow nilpotents and $R$-torsion in $S$.

The author's thesis and subsequent paper gave counterexamples to Koh's conjecture \cite{V93,Velez}. V\'elez and Fl\'orez later constructed lower-dimensional examples in characteristic two and in mixed characteristic \cite{VF}. Other positive results concern bigenerated extensions, bounded complexity, and annihilator and socle criteria \cite{GGV2016,VG2020}. Here we ask which curve singularities have $\K$, and which presentations force an individual extension to split in higher dimension.

\begin{maina}
Let $R$ be a one-dimensional Gorenstein local ring of residue characteristic $p>0$. If $R$ is weakly normal, it has $\K$. When $p=2$ and $2$ is zero or a nonzerodivisor in $R$, the converse holds. In equal characteristic $p$, weak normality is equivalent to $F$-purity.
\end{maina}

The point is that the values $\phi(1_S)$ of linear functionals form the conductor of an intermediate ring in the total quotient ring $Q(R)$. A determinant on a free resolution controls the powers of each scalar in that ring. Conversely, a missing square root yields a nonsplit extension over the given base, with minimal resolution $0\to R^6\to R^{10}\to S\to0$ (Theorem~\ref{realize}). The converse at odd primes remains undecided by this argument.

For an excellent one-dimensional local ring, the normalization of $(\widehat R)_{\mathrm{red}}$ is a product of complete discrete valuation rings. Say that its \emph{branch residue extensions are separable} if each of their residue fields is separable over $k$.

\begin{mainb}
An excellent seminormal local curve with separable branch residue extensions has $\K$. If $R$ is an excellent one-dimensional Cohen--Macaulay local ring with this separability property, residue characteristic two, and $2$ zero or a nonzerodivisor, then $\K$ is equivalent to seminormality. For an excellent Cohen--Macaulay local curve of equal characteristic $p>0$ with separable branch residue extensions, seminormality, weak normality and $F$-purity are equivalent.
\end{mainb}

Theorem B allows any number of branches and an imperfect residue field. Its proof starts with an exact criterion for a given extension of a fiber product of discrete valuation rings (Theorem~\ref{branch-split}). It uses the classical structure of seminormal and $F$-pure curves \cite{Swan,GT,GW,Gib,EH,HS}. Section~\ref{geometry} uses Cartier operators to pass from the generic point of a crossing to the whole singularity.

The last result in this introduction concerns cubic extensions. Let $I\subseteq\m$ and let $U$ be a $(b+2)\times b$ matrix with entries in $I$. On the closed fiber of the blow-up of $I$, its entries induce linear forms from $I/\m I$. We denote the resulting matrix by $\overline U_I$. Full column rank means rank $b$ at every geometric point of this fiber.

\begin{mainc}
Let $R\subseteq S$ be a module-finite extension with $S/R\simeq\operatorname{coker}U$, where $b\ge1$. Suppose $\overline U_I$ has full column rank and one of the following holds.
\begin{enumerate}[label=\textup{(\roman*)}]
\item $I$ is generated by an $R$-regular sequence of length three.
\item $I=\m$ and
$\operatorname{gr}_{\m}R\simeq k[X_0,X_1,X_2,X_3]/(f_2)$ for a nonzero quadratic form $f_2$.
\end{enumerate}
Then $R\subseteq S$ splits. In either case arbitrary changes of the entries in $\m I$ preserve the criterion.
\end{mainc}

There is no restriction on $b$, on the residue characteristic, or on the rank of $f_2$. For $I=(a_0,a_1,a_2)$ as in (i), multiplication of polynomials by $a_0+a_1Z+a_2Z^2$ gives an explicit admissible matrix for every $b$ (Corollary~\ref{toeplitz}). The proof uses the binary cubic of a rank-three algebra \cite[Theorem~2.1]{Wood}. In residue characteristic three its restriction to the exceptional fiber is controlled by a two-step filtration of a symmetric cube. For a hypersurface tangent cone of higher degree, Section~\ref{higher-surfaces} gives separate criteria for two and three relations.

In particular, Corollary~\ref{three-entries} handles all degrees and all numbers of relations when the entry ideal has three generators and the linear matrix has full rank on the exceptional surface.

Sections~\ref{homol} and \ref{ranks} treat individual extensions by means of $\Ext^1$, the tangent cone, and multiplicity. The quadratic result refines the argument of Hochster--McLaughlin \cite[Theorem~2.1]{HM}, replacing local factoriality by triviality of the reflexive determinant of the specified algebra. The target may have zero divisors.

\subsection*{Notation}
Write $M^*=\Hom_R(M,R)$, $\ell_R(M)$ for length, and
$H^0_{\m}(M)=\bigcup_n(0:_M\m^n)$. The \emph{evaluation ideal} of an extension is
\[
 J(S/R)=\{\phi(1_S):\phi\in\Hom_R(S,R)\}.
\]
It is the trace ideal, since $\phi(s)$ is the value at $1$ of $t\mapsto\phi(st)$. Thus $J(S/R)=R$ exactly when the extension splits. Evaluation commutes with flat base change by finite presentation, so a faithfully flat extension of the base detects splitting. Also, $1_S\notin\m S$ by Nakayama's lemma. If $S$ is free, its unit extends to a basis. In particular every local ring of depth zero has $\K$ by Auslander--Buchsbaum \cite{BH}.

A seminormal (respectively weakly normal) ring is reduced and admits no proper module-finite extension inside $Q(R)$ inducing a bijection on spectra and isomorphisms (respectively purely inseparable extensions) on residue fields. We use the square--cube criterion for seminormality and the root criterion of Lemma~\ref{wn}. In characteristic $p$, $F_*^eM$ denotes restriction of scalars along the $e$th Frobenius. The ring is $F$-pure when $R\to F_*R$ remains injective after tensoring with any module. We impose no finiteness condition on Frobenius.

One rank observation applies in every dimension. For a module-finite extension of finite projective dimension, lift multiplication by $s\in S$ to a chain map $f_\bullet(s)$ on a finite free resolution $P_\bullet$ of $S$. The alternating trace
\begin{equation}\label{perfect-trace}
 \tau(s)=\sum_i(-1)^i\operatorname{tr}(f_i(s))
 \quad\text{satisfies}\quad
 \tau(1)=r=\sum_i(-1)^i\operatorname{rank}P_i .
\end{equation}
Two lifts are homotopic. The alternating trace of $dh+hd$ vanishes by $\operatorname{tr}(AB)=\operatorname{tr}(BA)$, so $\tau$ is well defined and $R$-linear. At each minimal prime $S$ is free by Auslander--Buchsbaum, of rank $r$. Thus $r\in J(S/R)$ and $\tau/r$ is a retraction whenever $r$ is a unit. In particular all finite-projective-dimension extensions of generic rank three split if $3$ is a unit.
